\documentclass[a4paper,12pt]{report}
\usepackage[margin=0.5in]{geometry}
\usepackage{hyperref}
\usepackage{setspace}
\usepackage{lastpage}
\usepackage{tabularx}

\usepackage[nottoc,numbib]{tocbibind}
  % Preamble packages:
\usepackage{booktabs}
\usepackage{multirow}
\usepackage{caption}
\usepackage{multicol}
\usepackage{caption}



% for apa with apa
% \usepackage{natbib} 
% for ieee with IEEEtran
\usepackage[numbers]{natbib}
\usepackage{float}
\usepackage{grffile}  % Add to preamble
\usepackage{titlesec} % Add this line for chapter formatting

\usepackage{graphicx}

\usepackage{float} % If you use [H] placement


% Required in preamble
\usepackage{pgfplots}      % Core plotting package
\pgfplotsset{compat=1.18} % Set compatibility mode (use newest version available) % For general graphics handling (usually auto-loaded)

% Required packages in preamble
\usepackage{algorithm}      % For algorithm float environment
\usepackage{algpseudocode}  % For algorithmic structure (functions, loops, etc.)
\usepackage{xspace}         % For smart spacing after macros (optional but recommended)

\usepackage{csquotes}

\renewcommand{\familydefault}{\sfdefault}
\usepackage{helvet}
\usepackage{amsmath}
\usepackage{tikz}
\usetikzlibrary{positioning}

\renewcommand{\contentsname}{Table of Contents}

\usepackage{graphicx}
\usepackage{subcaption}
\usepackage{adjustbox}

\usepackage{booktabs} % For professional quality tables
\usepackage{caption}  % For captions of figures and tables
\usepackage{subcaption} % For subfigures within a figure environment
% \usepackage[backend=biber, style=apa]{biblatex} % Example for bibliography (if needed, adapt style)
% \addbibresource{references.bib} % Link to your bibliography file
% \usepackage{hyperref} % Optional: for clickable links and references
% \usepackage[capitalise]{cleveref} % Optional: for smart cross-referencing (e.g., Figure 1, Table 2)

\newenvironment{smallquote}
  {\begin{quote}\footnotesize\vspace{0em}}
  {\vspace{0em}\end{quote}}


\usepackage{pgffor}
% MAKE Smaller
% make lists small
\usepackage{enumitem}
\setlist{nosep} % Removes extra space in itemize/enumerate

\titlespacing*{\section}{0pt}{0.5em}{0.5em}
\titlespacing*{\subsection}{0pt}{0.5em}{0.5em}

% % \titleformat{\section}
% %   {\normalfont\normalsize\bfseries}
% %   {\thesection}{1em}{}

% \titleformat{\subsection}
%   {\normalfont\small\bfseries}
%   {\thesubsection}{1em}{}

\setlength{\parindent}{1.5em}
\setlength{\parskip}{0.5em} % Adjust spacing between paragraphs

% Customise hyperlink appearance (Optional)
\hypersetup{
    colorlinks=true,        % Enable colored links
    linkcolor=black,        % Color for internal links (TOC, sections, etc.)
    citecolor=black,        % Color for citation links
    filecolor=black,        % Color for file links
    urlcolor=blue          % Color for external URLs
}

% \titleformat{\chapter}[hang]{\normalfont\huge\bfseries}{\thechapter.}{0.5em}{}
% \titleformat{\chapter}[display]
%   {\normalfont\Large\bfseries}
%   {}
%   {0pt}
%   {\vspace{-2em}}
% \titlespacing*{\chapter}{0pt}{-2em}{1em}

% % Chapter title (smaller than default)
% \titleformat{\chapter}[display]
%   {\normalfont\Large\bfseries}
%   {} % No label
%   {0pt}
%   {\vspace{-2em}}

% % Section title (even smaller)
% \titleformat{\section}
%   {\normalfont\small\bfseries}
%   {\thesection}{1em}{}

% \titleformat{\chapter}[display]
%   {\normalfont\LARGE\bfseries}
%   {\thechapter}{1em}{} % <- this line sets the number before the title
% \titlespacing*{\chapter}{0pt}{-2em}{1em}

% \usepackage{titlesec}

\titleformat{\chapter}[block] % or [hang]
  {\normalfont\LARGE\bfseries}
  {\thechapter}{1em}{}

\titlespacing*{\chapter}{0pt}{-2em}{1em}


% Section: Large and bold
\titleformat{\section}
  {\normalfont\Large\bfseries}
  {\thesection}{1em}{}
\titlespacing*{\section}{0pt}{0.5em}{0.5em}

% Subsection: Slightly smaller than section, bold
\titleformat{\subsection}
  {\normalfont\large\bfseries}
  {\thesubsection}{1em}{}
\titlespacing*{\subsection}{0pt}{0.4em}{0.4em}

% Subsubsection: Normal size, bold (smaller than subsection)
\titleformat{\subsubsection}
  {\normalfont\normalsize\bfseries}
  {\thesubsubsection}{1em}{}
\titlespacing*{\subsubsection}{0pt}{0.3em}{0.3em}

\begin{document}


% LaTeX code for segmentation masks figure
% Requires: graphicx, subcaption, adjustbox packages

% Add these to your preamble if not already included:
% \usepackage{graphicx}
% \usepackage{subcaption}
% \usepackage{adjustbox}

% Option 2: Updated Paired Layout with Grad-CAM (Original → Mask → Grad-CAM triplets)
\begin{figure}[htbp]
    \centering
    \caption{Wing-tip segmentation and attention visualization pipeline for species identification. Each set shows the original image (left), its corresponding segmentation mask (middle) highlighting wing-tip patterns, and Grad-CAM visualization (right) showing model attention regions for morphological analysis.}
    \label{fig:segmentation_pairs}
    
    % Triplet 1
    \begin{subfigure}[b]{\textwidth}
        \centering
        \begin{subfigure}[b]{0.32\textwidth}
            \centering
            \includegraphics[width=\textwidth,height=3cm,keepaspectratio]{figures/original_1.png}
            \caption*{Original}
        \end{subfigure}%
        \hfill
        \begin{subfigure}[b]{0.32\textwidth}
            \centering
            \includegraphics[width=\textwidth,height=3cm,keepaspectratio]{figures/mask_1.png}
            \caption*{Segmentation}
        \end{subfigure}%
        \hfill
        \begin{subfigure}[b]{0.32\textwidth}
            \centering
            \includegraphics[width=\textwidth,height=3cm,keepaspectratio]{figures/gradcam_1.png}
            \caption*{Grad-CAM}
        \end{subfigure}
        \caption{Glaucous-winged Gull - Example 1}
    \end{subfigure}
    
    \vspace{0.5cm}
    
    % Triplet 2
    \begin{subfigure}[b]{\textwidth}
        \centering
        \begin{subfigure}[b]{0.32\textwidth}
            \centering
            \includegraphics[width=\textwidth,height=3cm,keepaspectratio]{figures/original_2.png}
            \caption*{Original}
        \end{subfigure}%
        \hfill
        \begin{subfigure}[b]{0.32\textwidth}
            \centering
            \includegraphics[width=\textwidth,height=3cm,keepaspectratio]{figures/mask_2.png}
            \caption*{Segmentation}
        \end{subfigure}%
        \hfill
        \begin{subfigure}[b]{0.32\textwidth}
            \centering
            \includegraphics[width=\textwidth,height=3cm,keepaspectratio]{figures/gradcam_2.png}
            \caption*{Grad-CAM}
        \end{subfigure}
        \caption{Glaucous-winged Gull - Example 2}
    \end{subfigure}
    
    \vspace{0.5cm}
    
    % % Triplet 3
    % \begin{subfigure}[b]{\textwidth}
    %     \centering
    %     \begin{subfigure}[b]{0.32\textwidth}
    %         \centering
    %         \includegraphics[width=\textwidth,height=3cm,keepaspectratio]{figures/original_3.png}
    %         \caption*{Original}
    %     \end{subfigure}%
    %     \hfill
    %     \begin{subfigure}[b]{0.32\textwidth}
    %         \centering
    %         \includegraphics[width=\textwidth,height=3cm,keepaspectratio]{figures/mask_3.png}
    %         \caption*{Segmentation}
    %     \end{subfigure}%
    %     \hfill
    %     \begin{subfigure}[b]{0.32\textwidth}
    %         \centering
    %         \includegraphics[width=\textwidth,height=3cm,keepaspectratio]{figures/gradcam_3.png}
    %         \caption*{Grad-CAM}
    %     \end{subfigure}
    %     \caption{Slaty-backed Gull - Example 1}
    % \end{subfigure}
    
    % \vspace{0.5cm}
    
    % % Triplet 4
    % \begin{subfigure}[b]{\textwidth}
    %     \centering
    %     \begin{subfigure}[b]{0.32\textwidth}
    %         \centering
    %         \includegraphics[width=\textwidth,height=3cm,keepaspectratio]{figures/original_4.png}
    %         \caption*{Original}
    %     \end{subfigure}%
    %     \hfill
    %     \begin{subfigure}[b]{0.32\textwidth}
    %         \centering
    %         \includegraphics[width=\textwidth,height=3cm,keepaspectratio]{figures/mask_4.png}
    %         \caption*{Segmentation}
    %     \end{subfigure}%
    %     \hfill
    %     \begin{subfigure}[b]{0.32\textwidth}
    %         \centering
    %         \includegraphics[width=\textwidth,height=3cm,keepaspectratio]{figures/gradcam_4.png}
    %         \caption*{Grad-CAM}
    %     \end{subfigure}
    %     \caption{Slaty-backed Gull - Example 2}
    % \end{subfigure}
\end{figure}

% Option 1: Basic 2x4 Layout (Recommended for publication)
\begin{figure}[htbp]
    \centering
    \caption{Image segmentation pipeline showing original gull images and corresponding wing-tip segmentation masks. (A) Original images of Glaucous-winged Gull (GWG) and Slaty-backed Gull (SBG) specimens. (B) Corresponding segmentation masks highlighting wing-tip regions for morphological analysis.}
    \label{fig:segmentation_pipeline}
    
    % Top row - Original images
    \begin{subfigure}[b]{0.24\textwidth}
        \centering
        \includegraphics[width=\textwidth,height=3cm,keepaspectratio]{figures/original_1.png}
        \caption*{GWG - Example 1}
    \end{subfigure}%
    \hfill
    \begin{subfigure}[b]{0.24\textwidth}
        \centering
        \includegraphics[width=\textwidth,height=3cm,keepaspectratio]{figures/original_2.png}
        \caption*{GWG - Example 2}
    \end{subfigure}%
    \hfill
    \begin{subfigure}[b]{0.24\textwidth}
        \centering
        \includegraphics[width=\textwidth,height=3cm,keepaspectratio]{figures/original_3.png}
        \caption*{SBG - Example 1}
    \end{subfigure}%
    \hfill
    \begin{subfigure}[b]{0.24\textwidth}
        \centering
        \includegraphics[width=\textwidth,height=3cm,keepaspectratio]{figures/original_4.png}
        \caption*{SBG - Example 2}
    \end{subfigure}
    
    \vspace{0.5cm}
    \textbf{(A) Original Images}
    \vspace{0.3cm}
    
    % Bottom row - Segmentation masks
    \begin{subfigure}[b]{0.24\textwidth}
        \centering
        \includegraphics[width=\textwidth,height=3cm,keepaspectratio]{figures/mask_1.png}
        \caption*{Mask 1}
    \end{subfigure}%
    \hfill
    \begin{subfigure}[b]{0.24\textwidth}
        \centering
        \includegraphics[width=\textwidth,height=3cm,keepaspectratio]{figures/mask_2.png}
        \caption*{Mask 2}
    \end{subfigure}%
    \hfill
    \begin{subfigure}[b]{0.24\textwidth}
        \centering
        \includegraphics[width=\textwidth,height=3cm,keepaspectratio]{figures/mask_3.png}
        \caption*{Mask 3}
    \end{subfigure}%
    \hfill
    \begin{subfigure}[b]{0.24\textwidth}
        \centering
        \includegraphics[width=\textwidth,height=3cm,keepaspectratio]{figures/mask_4.png}
        \caption*{Mask 4}
    \end{subfigure}
    
    \vspace{0.2cm}
    \textbf{(B) Segmentation Masks}
\end{figure}

% Option 2: Paired Layout (Original → Mask pairs)
\begin{figure}[htbp]
    \centering
    \caption{Wing-tip segmentation Examples demonstrating the image processing pipeline for species identification. Each pair shows the original image (left) and its corresponding segmentation mask (right) highlighting wing-tip patterns used for morphological analysis.}
    \label{fig:segmentation_pairs}
    
    % Pair 1
    \begin{subfigure}[b]{0.48\textwidth}
        \centering
        \begin{subfigure}[b]{0.48\textwidth}
            \centering
            \includegraphics[width=\textwidth,height=3cm,keepaspectratio]{figures/original_1.png}
            \caption*{Original}
        \end{subfigure}%
        \hfill
        \begin{subfigure}[b]{0.48\textwidth}
            \centering
            \includegraphics[width=\textwidth,height=3cm,keepaspectratio]{figures/mask_1.png}
            \caption*{Segmentation}
        \end{subfigure}
        \caption{Glaucous-winged Gull - Example 1}
    \end{subfigure}%
    \hfill
    % Pair 2
    \begin{subfigure}[b]{0.48\textwidth}
        \centering
        \begin{subfigure}[b]{0.48\textwidth}
            \centering
            \includegraphics[width=\textwidth,height=3cm,keepaspectratio]{figures/original_2.png}
            \caption*{Original}
        \end{subfigure}%
        \hfill
        \begin{subfigure}[b]{0.48\textwidth}
            \centering
            \includegraphics[width=\textwidth,height=3cm,keepaspectratio]{figures/mask_2.png}
            \caption*{Segmentation}
        \end{subfigure}
        \caption{Glaucous-winged Gull - Example 2}
    \end{subfigure}
    
    \vspace{0.5cm}
    
    % Pair 3
    \begin{subfigure}[b]{0.48\textwidth}
        \centering
        \begin{subfigure}[b]{0.48\textwidth}
            \centering
            \includegraphics[width=\textwidth,height=3cm,keepaspectratio]{figures/original_3.png}
            \caption*{Original}
        \end{subfigure}%
        \hfill
        \begin{subfigure}[b]{0.48\textwidth}
            \centering
            \includegraphics[width=\textwidth,height=3cm,keepaspectratio]{figures/mask_3.png}
            \caption*{Segmentation}
        \end{subfigure}
        \caption{Slaty-backed Gull - Example 1}
    \end{subfigure}%
    \hfill
    % Pair 4
    \begin{subfigure}[b]{0.48\textwidth}
        \centering
        \begin{subfigure}[b]{0.48\textwidth}
            \centering
            \includegraphics[width=\textwidth,height=3cm,keepaspectratio]{figures/original_4.png}
            \caption*{Original}
        \end{subfigure}%
        \hfill
        \begin{subfigure}[b]{0.48\textwidth}
            \centering
            \includegraphics[width=\textwidth,height=3cm,keepaspectratio]{figures/mask_4.png}
            \caption*{Segmentation}
        \end{subfigure}
        \caption{Slaty-backed Gull - Example 2}
    \end{subfigure}
\end{figure}

% Option 3: Compact Grid Layout
\begin{figure}[htbp]
    \centering
    \caption{Segmentation pipeline results for wing-tip pattern analysis. Top row: original images; bottom row: corresponding segmentation masks highlighting wing-tip regions for morphological feature extraction.}
    \label{fig:segmentation_grid}
    
    % Create a table-like structure
    \begin{tabular}{cccc}
        \includegraphics[width=0.22\textwidth,height=3cm,keepaspectratio]{figures/original_1.png} &
        \includegraphics[width=0.22\textwidth,height=3cm,keepaspectratio]{figures/original_2.png} &
        \includegraphics[width=0.22\textwidth,height=3cm,keepaspectratio]{figures/original_3.png} &
        \includegraphics[width=0.22\textwidth,height=3cm,keepaspectratio]{figures/original_4.png} \\
        \small{GWG-1 Original} & \small{GWG-2 Original} & \small{SBG-1 Original} & \small{SBG-2 Original} \\[0.3cm]
        
        \includegraphics[width=0.22\textwidth,height=3cm,keepaspectratio]{figures/mask_1.png} &
        \includegraphics[width=0.22\textwidth,height=3cm,keepaspectratio]{figures/mask_2.png} &
        \includegraphics[width=0.22\textwidth,height=3cm,keepaspectratio]{figures/mask_3.png} &
        \includegraphics[width=0.22\textwidth,height=3cm,keepaspectratio]{figures/mask_4.png} \\
        \small{GWG-1 Mask} & \small{GWG-2 Mask} & \small{SBG-1 Mask} & \small{SBG-2 Mask} \\
    \end{tabular}
\end{figure}

% Option 4: Publication-ready with detailed captions
\begin{figure}[!htbp]
    \centering
    \caption{Image segmentation pipeline for automated wing-tip pattern analysis in \textit{Larus} species. Representative Examples show the segmentation process applied to field photographs of Glaucous-winged Gull (\textit{L. glaucescens}) and Slaty-backed Gull (\textit{L. schistisagus}). (A) Original field photographs showing wing-tip regions. (B) Corresponding binary segmentation masks generated through automated image processing, highlighting wing-tip patterns used for morphological feature extraction and species classification.}
    \label{fig:segmentation_detailed}
    
    % Row A - Original images
    \noindent\textbf{(A) Original Images}\\[0.2cm]
    \begin{subfigure}[t]{0.24\textwidth}
        \centering
        \adjustbox{frame=0.5pt}{\includegraphics[width=\textwidth,height=3.5cm,keepaspectratio]{figures/original_1.png}}
        \caption*{\small GWG specimen 1}
        \label{fig:orig1}
    \end{subfigure}%
    \hfill
    \begin{subfigure}[t]{0.24\textwidth}
        \centering
        \adjustbox{frame=0.5pt}{\includegraphics[width=\textwidth,height=3.5cm,keepaspectratio]{figures/original_2.png}}
        \caption*{\small GWG specimen 2}
        \label{fig:orig2}
    \end{subfigure}%
    \hfill
    \begin{subfigure}[t]{0.24\textwidth}
        \centering
        \adjustbox{frame=0.5pt}{\includegraphics[width=\textwidth,height=3.5cm,keepaspectratio]{figures/original_3.png}}
        \caption*{\small SBG specimen 1}
        \label{fig:orig3}
    \end{subfigure}%
    \hfill
    \begin{subfigure}[t]{0.24\textwidth}
        \centering
        \adjustbox{frame=0.5pt}{\includegraphics[width=\textwidth,height=3.5cm,keepaspectratio]{figures/original_4.png}}
        \caption*{\small SBG specimen 2}
        \label{fig:orig4}
    \end{subfigure}
    
    \vspace{0.8cm}
    
    % Row B - Segmentation masks
    \noindent\textbf{(B) Segmentation Masks}\\[0.2cm]
    \begin{subfigure}[t]{0.24\textwidth}
        \centering
        \adjustbox{frame=0.5pt}{\includegraphics[width=\textwidth,height=3.5cm,keepaspectratio]{figures/mask_1.png}}
        \caption*{\small Wing-tip mask 1}
        \label{fig:mask1}
    \end{subfigure}%
    \hfill
    \begin{subfigure}[t]{0.24\textwidth}
        \centering
        \adjustbox{frame=0.5pt}{\includegraphics[width=\textwidth,height=3.5cm,keepaspectratio]{figures/mask_2.png}}
        \caption*{\small Wing-tip mask 2}
        \label{fig:mask2}
    \end{subfigure}%
    \hfill
    \begin{subfigure}[t]{0.24\textwidth}
        \centering
        \adjustbox{frame=0.5pt}{\includegraphics[width=\textwidth,height=3.5cm,keepaspectratio]{figures/mask_3.png}}
        \caption*{\small Wing-tip mask 3}
        \label{fig:mask3}
    \end{subfigure}%
    \hfill
    \begin{subfigure}[t]{0.24\textwidth}
        \centering
        \adjustbox{frame=0.5pt}{\includegraphics[width=\textwidth,height=3.5cm,keepaspectratio]{figures/mask_4.png}}
        \caption*{\small Wing-tip mask 4}
        \label{fig:mask4}
    \end{subfigure}
\end{figure}

% Required packages for preamble:
% \usepackage{graphicx}
% \usepackage{subcaption}
% \usepackage{adjustbox}
% \usepackage{array}

% Folder structure needed:
% your_project/
% └── figures/
%     ├── original_1.png  (GWG Example 1)
%     ├── original_2.png  (GWG Example 2)
%     ├── original_3.png  (SBG Example 1)
%     ├── original_4.png  (SBG Example 2)
%     ├── mask_1.png      (mask for original_1)
%     ├── mask_2.png      (mask for original_2)
%     ├── mask_3.png      (mask for original_3)
%     └── mask_4.png      (mask for original_4)

% Usage Instructions:
% 1. Add the required packages to your preamble
% 2. Create a 'figures' folder in your LaTeX project directory
% 3. Place your 8 images (4 originals + 4 masks) in the figures folder
% 4. Rename them according to the naming convention above
% 5. Choose one of the four options above and add it to your document
% 6. Compile your LaTeX document

% Tips:
% - Option 1: Best for general publication use
% - Option 2: Good for showing clear original→mask relationships
% - Option 3: Most compact, good for space-constrained journals
% - Option 4: Most detailed, good for method papers

% Customization options:
% - Adjust width/height parameters to fit your layout
% - Change frame thickness by modifying frame=0.5pt
% - Modify captions and labels as needed
% - Add/remove species labels as appropriate

\end{document}